%% kapa: exam=1 round=36 year=2025 period=1 form=A subject=real_estate subject_ko=부동산학원론 seq=40 number=120
%% source: Q-Net 감정평가사 기출문제 (cst00309) artlSeq=5231525 file=raw/1008/5231525/2025년 제36회 감정평가사 1차시험 1교시 문제지 원본.pdf
%% conversion: kordoc 4.24.51 -> tools/gen1.py (auto); review=proofread-sonnet; flags=none
\begin{question}{120}
감정평가에 관한 규칙상 ‘적정한 실거래가’에 관한 설명으로 옳은 것은?
\begin{choices}
  \item 도시지역의 경우, 거래시점이 4년 이내의 것이어야 한다.
  \item 도시지역이 아닌 경우, 거래시점이 6년 이내의 것이어야 한다.
  \item 적정한 실거래가의 기준이 되는 도시지역에 계획관리지역이 포함된다.
  \item ｢부동산 거래신고 등에 관한 법률｣에 따라 신고된 실제 거래가격이어야 한다.
  \item 실거래가는 인근지역 지가수준과의 차이와 관계없이 적정한 실거래가로 인정되어야 한다.
\end{choices}
\end{question}
\answer{4}
